% =============================================================================
% MethodAtlas Documentation — LaTeX preamble extension
% Included via pandoc -H; appended to the auto-generated preamble.
% Compatible with: XeLaTeX + MikTeX (Windows) / TeX Live (Linux / macOS)
% =============================================================================

% ── XeLaTeX logo ─────────────────────────────────────────────────────────────
% \XeLaTeX is normally provided by xltxtra, which pandoc does not load.
% Define it here as a fallback so back-matter.tex can use \XeLaTeX safely.
\providecommand{\XeLaTeX}{Xe\LaTeX}

% ── Colour palette ────────────────────────────────────────────────────────────
% xcolor is already loaded by pandoc's template; we only define our colours.
\definecolor{MA@navy}{HTML}{1B2A4A}
\definecolor{MA@teal}{HTML}{00848E}
\definecolor{MA@codebg}{HTML}{F7F7F9}
\definecolor{MA@codeframe}{HTML}{DADCE0}
\definecolor{MA@comment}{HTML}{5C7A29}
\definecolor{MA@string}{HTML}{C0392B}

% Override pandoc's hyperref colours with our palette.
% Deferred to \AtBeginDocument so hyperref is guaranteed to be loaded first.
% (pandoc inserts -H content before its own \usepackage{hyperref}.)
\AtBeginDocument{%
  \hypersetup{%
    linkcolor=MA@navy,%
    urlcolor=MA@teal,%
    citecolor=MA@navy,%
    filecolor=MA@navy,%
    pdftitle={MethodAtlas Documentation},%
    pdfauthor={Egothor / Accenture},%
    pdfsubject={Security test discovery, AI classification, SARIF 2.1.0},%
    pdfkeywords={MethodAtlas, security testing, test inventory, SARIF, AI, OWASP},%
  }%
}

% ── Unicode character fallbacks ───────────────────────────────────────────────
% Map Unicode characters that Latin Modern lacks to equivalent LaTeX commands.
% newunicodechar intercepts at the TeX level, before the font engine, so it
% works correctly with XeLaTeX + fontspec.
\usepackage{newunicodechar}
\newunicodechar{≥}{\ensuremath{\geq}}        % U+2265 GREATER-THAN OR EQUAL TO
\newunicodechar{↪}{\ensuremath{\hookrightarrow}} % U+21AA RIGHTWARDS ARROW WITH HOOK
\newunicodechar{⚠}{\textbf{[!]}}             % U+26A0 WARNING SIGN
% U+FE0F VARIATION SELECTOR-16 is an invisible emoji modifier; set its
% catcode to 9 (ignored) so XeLaTeX silently drops it everywhere.
\catcode`\^^^^fe0f=9

% ── Subliminal refinements ────────────────────────────────────────────────────
\usepackage{microtype}   % kerning, tracking, protrusion — perceptible improvement

% ── Code-block line wrapping (key feature) ────────────────────────────────────
% fvextra extends fancyvrb with breaklines and breakanywhere.
% pandoc wraps every highlighted code block in the Highlighting environment
% (a subclass of Verbatim).  Redefining it here with breaklines=true causes
% long command lines, URLs, and deeply nested JSON/YAML to wrap gracefully
% instead of running off the right margin.
\usepackage{fvextra}
\DefineVerbatimEnvironment{Highlighting}{Verbatim}{%
  breaklines=true,%
  breakanywhere=true,%
  breakanywheresymbolpre={},%
  breaksymbol={\mbox{\tiny$\hookrightarrow$}},%
  breaksymbolright={},%
  commandchars=\\\{\}%
}

% ── Better tables ─────────────────────────────────────────────────────────────
\usepackage{booktabs}    % \toprule, \midrule, \bottomrule
\usepackage{longtable}   % tables that span page breaks (pandoc emits these)
\usepackage{array}       % extended column specifiers

% ── Section heading typography ────────────────────────────────────────────────
\usepackage{titlesec}

% Chapter: large, navy, ruled below
\titleformat{\chapter}[display]
  {\normalfont\LARGE\bfseries\color{MA@navy}}
  {\chaptertitlename~\thechapter}
  {10pt}
  {\huge}
  [\vspace{2pt}{\color{MA@teal}\titlerule[1.2pt]}]

% Section: bold navy
\titleformat{\section}
  {\normalfont\Large\bfseries\color{MA@navy}}
  {\thesection}{1em}{}

% Subsection: medium weight, teal
\titleformat{\subsection}
  {\normalfont\large\bfseries\color{MA@teal}}
  {\thesubsection}{1em}{}

% Subsubsection: normal weight, navy
\titleformat{\subsubsection}
  {\normalfont\normalsize\bfseries\color{MA@navy}}
  {\thesubsubsection}{1em}{}

% ── Running headers and footers ───────────────────────────────────────────────
\usepackage{fancyhdr}
\usepackage{emptypage}   % suppress headers/footers on empty pages

\pagestyle{fancy}
\fancyhf{}
\fancyhead[LE]{\small\color{MA@navy}\leftmark}        % left even: chapter title
\fancyhead[RO]{\small\color{MA@navy}\rightmark}       % right odd: section title
\fancyfoot[LE,RO]{\small\thepage}
\fancyfoot[C]{\small\color{MA@teal}MethodAtlas Documentation}
\renewcommand{\headrulewidth}{0.4pt}
\renewcommand{\footrulewidth}{0pt}

% Plain style for chapter opening pages
\fancypagestyle{plain}{%
  \fancyhf{}%
  \fancyfoot[LE,RO]{\small\thepage}%
  \renewcommand{\headrulewidth}{0pt}%
}

% ── Custom title page ─────────────────────────────────────────────────────────
% Redefine \maketitle to produce a full-bleed cover page in navy.
% pandoc calls \maketitle at the start of the document body when a title is set.
% \makeatletter is required so that \@date (an internal LaTeX macro, with @
% as part of the control-sequence name) is accessible in user code loaded
% via pandoc -H, where @ has catcode 12 (punctuation) by default.
\makeatletter
\renewcommand{\maketitle}{%
  \begin{titlepage}%
    \pagecolor{MA@navy}%
    \color{white}%
    \vspace*{3cm}%
    \begin{center}%
      % Logotype: Method in white, Atlas in teal
      {\fontsize{54}{64}\selectfont\bfseries%
        Method\textcolor{MA@teal}{Atlas}}\\[0.6em]
      {\large\itshape Security Test Intelligence Platform}\\[3.5em]
      {\color{MA@teal}\rule{0.42\textwidth}{0.8pt}}\\[2.5em]
      {\Large\bfseries Developer Reference}\\[0.7em]
      {\normalsize Full Documentation Set}\\[4.5em]
      {\small Security test discovery\enspace·\enspace AI classification\enspace·\enspace SARIF\,2.1.0}\\[0.35em]
      {\small OWASP SAMM\enspace·\enspace NIST SSDF\enspace·\enspace
              ISO\,27001\enspace·\enspace DORA\enspace·\enspace PCI\,DSS\enspace·\enspace SOC\,2}%
    \end{center}%
    \vfill%
    \begin{center}%
      {\small\color{white!80}\@date}\\[0.5em]
      {\footnotesize\textcopyright\enspace Egothor\,/\,Accenture
        \quad\textbullet\quad
        Apache\,License\,2.0}\\[0.4em]
      {\footnotesize ISBN:\enspace(not yet assigned)}%
    \end{center}%
    \vspace{1.5cm}%
  \end{titlepage}%
  % Verso of title page — copyright details on white
  \nopagecolor%
  \thispagestyle{empty}%
  \vspace*{\fill}%
  \noindent\small%
  \textbf{MethodAtlas — Developer Reference}\\[0.4em]
  Copyright \textcopyright\ 2024--\the\year\ Egothor s.r.o.\ and Accenture.\\[0.3em]
  Licensed under the Apache License, Version 2.0.  You may obtain a copy of the
  licence at \url{https://www.apache.org/licenses/LICENSE-2.0}.\\[0.3em]
  Unless required by applicable law or agreed to in writing, software distributed
  under the Licence is distributed on an \textsc{``as is''} basis, without
  warranties or conditions of any kind, either express or implied.\\[0.6em]
  Source code and current online documentation:\\
  \url{https://accenture.github.io/MethodAtlas/}\\[0.6em]
  \textit{ISBN: (not yet assigned)}\\[0.4em]
  \textit{First edition.  Built: \@date.}%
  \vspace*{1cm}%
  \clearpage%
}
\makeatother
